\documentclass[a4paper,11pt]{article}
\usepackage[utf8]{inputenc}
\usepackage[T1]{fontenc}
\usepackage[french]{babel}
\usepackage{lmodern}
\usepackage{amsmath,amssymb,amsthm}
\usepackage{mathrsfs}
\usepackage{enumitem}
\usepackage{multicol}

\setlist[enumerate]{label=\textbf{\textcolor{Couleur8}{\arabic*.}}, left=1em}

% Si vous voulez aussi modifier les listes enumerate imbriquées :
\setlist[enumerate,2]{label=\textcolor{Couleur8}{\alph*)}}
\setlist[enumerate,3]{label=\textcolor{Couleur8}{\roman*)}}
\setlist[enumerate,4]{label=\textcolor{Couleur8}{\Alph*)}}
\usepackage{graphicx}
\usepackage{xcolor}
\usepackage{tcolorbox}
\usepackage{geometry}
\usepackage{fancyhdr}
\usepackage{titlesec}
\usepackage{cancel}
\usepackage{ifthen}
\usepackage{tikz}
\usetikzlibrary{arrows.meta}
\usepackage{tkz-tab}
\usepackage{eurosym}

% OPTION PROF/ÉLÈVE
% Décommentez UNE des deux lignes suivantes :
\newboolean{prof}\setboolean{prof}{true}   % Version professeur (avec corrections des devoirs)
%\newboolean{prof}\setboolean{prof}{false}  % Version élève (sans corrections des devoirs)

\definecolor{Couleur1}{HTML}{4A5D7A} %Th
\definecolor{Couleur2}{HTML}{A5B5D6} %Prop
\definecolor{Couleur3}{HTML}{A8C68A} %Méthode
\definecolor{Couleur6}{HTML}{8A9CC1} %Def
\definecolor{Couleur7}{HTML}{E6B459} %Rq Voc
\definecolor{Couleur8}{HTML}{D36740} %Titres
\definecolor{correction}{HTML}{D36740} % Couleur pour les corrections
\definecolor{coopmathscorr}{HTML}{F15929} % Couleur corrections coopmaths

% Configuration des marges
\geometry{
	top=2cm,
	bottom=2cm,
	inner=1.5cm,
	outer=1.5cm,
}

% Police
\renewcommand{\familydefault}{\sfdefault}

% Compteurs
\newcounter{seancenum}
\newcounter{seancenumD}
\setcounter{seancenum}{1}
\setcounter{seancenumD}{1}

% Configuration des en-têtes et pieds de page
\pagestyle{fancy}
\fancyhf{}
\ifthenelse{\boolean{prof}}{
    \fancyhead[L]{\textcolor{Couleur8}{\textbf{Corrections Automatismes - Classe de seconde - Version Professeur}}}
}{
    \fancyhead[L]{\textcolor{Couleur8}{\textbf{Corrections Automatismes - Séance 20 - Classe de seconde}}}
}
\fancyhead[R]{\textcolor{Couleur8}{\textbf{2025-2026}}}
\fancyfoot[L]{\textcolor{Couleur8}{Mme Bertail et Mme Pinto}}
\fancyfoot[R]{\textcolor{Couleur8}{\thepage}}
\renewcommand{\headrulewidth}{1pt}
\renewcommand{\headrule}{\hbox to\headwidth{\color{Couleur8}\leaders\hrule height \headrulewidth\hfill}}

% Environnements personnalisés
\newtcolorbox{seance}[1]{
    colback=Couleur6!10,
    colframe=Couleur1!65,
    fonttitle=\bfseries\large,
    title=Correction Séance \theseancenum #1,
    left=5mm,
    right=5mm,
    top=3mm,
    bottom=3mm,
    arc=0mm,
    after=\stepcounter{seancenum}
}

\newtcolorbox{seanceD}[1]{
    colback=Couleur6!5,
    colframe=Couleur6!45,
    fonttitle=\bfseries\large,
    title=Correction Devoirs - Séance \theseancenumD #1,
    left=5mm,
    right=5mm,
    top=3mm,
    bottom=3mm,
    arc=0mm,
    after=\stepcounter{seancenumD}
}

\begin{document}

\setcounter{seancenum}{20}
\newpage
\begin{seance}{} %20

\begin{enumerate}
\item On cherche s'il existe $x\in \mathbb{R}$ tel que $f(x)=54$.\\
On résout cette équation en isolant le carré.\\
$\begin{aligned}
8x^2-2&=54\\
8x^2&=56\\
x^2&=\dfrac{56}{8}\\
x^2&=7
\end{aligned}$\\
Cette équation a deux solutions : $-\sqrt{7}$ et $\sqrt{7}$.\\
Les antécédents de $54$ sont : ${\color{correction}\boldsymbol{-\sqrt{7}}}$ et ${\color{correction}\boldsymbol{\sqrt{7}}}$.

\item On arrondit $590$ à $600$ et $379$ à $400$.\\
On obtient : $600 \times 400 = 240\,000$.\\
Un ordre de grandeur de $590 \times 379$ est donc ${\color{correction}\boldsymbol{240\,000}}$.\\
La bonne réponse est la réponse {\bfseries \color{correction}B}.

\item Dans $1$ heure, il y a $4\times 15$ minutes.\\
Ainsi, en $1$ heure, elle parcourt $15$ fois plus de distance, soit $0{,}4\times 15=6\text{ km}$.\\
Sa vitesse moyenne est donc ${\color{correction}\boldsymbol{6}}\text{ km/h}$.

\item $\dfrac{16}{49}x^2+\dfrac{40}{7}x+25=\left(\dfrac{4}{7}x\right)^2+2 \times \dfrac{4}{7}x \times 5 + 5^2={\color{correction}\boldsymbol{\left(\dfrac{4}{7}x+5\right)^2}}$

\item Pour résoudre graphiquement cette inéquation :\\
$\bullet$ On trace la parabole d'équation $y=x^2$.\\
$\bullet$ On trace la droite horizontale d'équation $y=12$. Cette droite coupe la parabole en $-\sqrt{12}$ et $\sqrt{12}$.\\
$\bullet$ Les solutions de l'inéquation sont les abscisses des points de la courbe qui se situent sur ou au-dessus de la droite.\\
On en déduit que l'ensemble des solutions de l'inéquation est : {\bfseries \color{correction}$S = ]-\infty\,;\,-\sqrt{12}] \cup [\sqrt{12}\,;\,+\infty[$}.

\end{enumerate}

\end{seance}

\end{document}